% ============================================================================
\chapter{HTTP en PHP~: requ\^etes, formulaires, sessions, authentification}
\label{ch-php-http}
% ============================================================================

\begin{objectif}
\begin{itemize}
  \item Lire une requ\^ete proprement~: m\'ethodes, en-t\^etes, \verb|$_GET|, \verb|$\_POST|, corps JSON.
  \item Envoyer une r\'eponse correcte~: codes, en-t\^etes, redirections, streaming.
  \item Construire un formulaire robuste~: valeurs conserv\'ees, erreurs, protection CSRF.
  \item G\'erer une session et un mot de passe de bout en bout~: connexion, d\'econnexion,
        \og ~mot de passe oubli\'e~ \fg{}, r\'eg\'en\'eration d'identifiant.
  \item Uploader un fichier sans ouvrir de faille, et envoyer un email.
\end{itemize}
\end{objectif}

\section{Requ\^ete et r\'eponse~: le contrat}

\begin{codephp}
<?php
declare(strict_types=1);

// ---- LIRE LA REQUETE
$modele = $_SERVER['REQUEST_METHOD'];                       // 'GET' | 'POST' | 'PUT' ...
$uri = parse_url($_SERVER['REQUEST_URI'] ?? '/', PHP_URL_PATH);
$query = $uri;
$entete = fn (string $nom): ?string => $_SERVER['HTTP_' . strtoupper(str_replace('-', '_', $nom))] ?? null;
$contentType = $entete('content-type') ?? '';
$origine = $entete('origin');
$ip = $_SERVER['REMOTE_ADDR'] ?? 'inconnue';              // avec un proxy : X-Forwarded-For, a verifier
$estAjax = str_contains((string) $entete('x-requested-with'), 'XMLHttpRequest');

// Les superglobales, et ce qu'il faut en savoir
$_GET;              // chaine uniquement ; attention : tableau associatif brut, non filtre
$_POST;             // idem pour application/x-www-form-urlencoded et multipart/form-data
$_FILES;            // upload (voir plus bas)
$_COOKIE;           // cookies lus
$_SESSION;          // voir session
$_SERVER;           // meta : METHODE, URI, HTTP_*, DOCUMENT_ROOT, REQUEST_TIME_FLOAT
$_REQUEST;          // MELANGE des trois precedents selon variables_order : NE L'UTILISE PAS
$_ENV;              // seulement si variables_order contient "G" ; sinon : getenv() ou les constants

// Corps JSON (les API modernes ne passent plus par $_POST)
$brut = file_get_contents('php://input');
$json = $contentType !== '' && str_contains($contentType, 'json')
    ? json_decode($brut ?: 'null', true, 512, JSON_THROW_ON_ERROR)
    : null;
$entrees = $json ?? $_POST;

// Lecture securisee + validation en un point : mon petit "getter" de controleur
function champ(array $source, string $cle, ?string $defaut = null): ?string
{
    $v = $source[$cle] ?? $defaut;
    return is_string($v) ? trim($v) : ($v === null ? null : (string) $v);
}
function champInt(array $source, string $cle, ?int $defaut = null): ?int
{
    $v = $source[$cle] ?? null;
    return is_numeric($v) ? (int) $v : $defaut;
}
function champBool(array $source, string $cle): bool
{
    return in_array($source[$cle] ?? null, ['1', 1, 'on', 'true', true], true);
}
function filtre(array $source, string $cle, int $filtrePhp, array $options = []): mixed
{
    return filter_input_array($source === $_GET ? INPUT_GET : INPUT_POST, [$cle => $options ?: $filtrePhp])[$cle] ?? false;
}
$page = max(1, champInt($_GET, 'page', 1));
$email = filter_var(champ($entrees, 'email') ?? '', FILTER_VALIDATE_EMAIL);   // false si invalide
$urls = filter_var($entrees['site'] ?? '', FILTER_VALIDATE_URL);
$id = filter_var($_GET['id'] ?? null, FILTER_VALIDATE_INT, ['options' => ['min_range' => 1]]);

// ---- ENVOYER LA REPONSE
http_response_code(201);
header('content-type: application/json; charset=utf-8');
header('cache-control: private, no-store');
header('x-page: ' . $page);
header_link: header('link: </style.css>; rel=preload; as=style', false);   // false = ne pas remplacer
header('retry-after: 60');
echo json_encode(['ok' => true, 'page' => $page], JSON_UNESCAPED_UNICODE | JSON_THROW_ON_ERROR);
exit;                                            // apres une redirection : TOUJOURS exit

// Redirections : 303 apres un POST (modele POST/REDIRECT/GET), 301 si c'est definitif
function rediriger(string $url, int $code = 303): never
{
    header('location: ' . $url, true, $code);
    exit;
}
function telecharger(string $chemin, string $nom): never
{
    header('content-type: text/csv; charset=utf-8');
    header('content-disposition: attachment; filename="' . addslashes($nom) . '"');
    header('content-length: ' . (string) filesize($chemin));
    readfile($chemin);
    exit;
}
\end{codephp}

\begin{piege}
\verb|header()| \'echoue avec \og ~headers already sent~ \fg{} d\`es qu'un octet a \'et\'e \'emis~: un espace
avant \verb|<?php|, un saut de ligne apr\`es un \verb|?>|, un \verb|echo| de debug. Les
r\`egles qui \'eliminent ce bug pour toujours~: fichiers PHP \emph{sans} balise fermante finale,
\verb|ob\_start()| dans le bootstrap si tu dois bufferiser, et aucune sortie avant la logique.
\end{piege}

\section{Un formulaire qui se tient~: le motif complet}

\begin{codephp}
<?php
// src/Action/Inscription.php  ---  un seul fichier, trois cas : affichage, rejet, succes.
declare(strict_types=1);

session_start();
$csrf = $_SESSION['csrf'] ??= bin2hex(random_bytes(32));
$valeurs = ['pseudo' => '', 'email' => '', 'message' => '', 'newsletter' => '1'];
$erreurs = [];

if ($_SERVER['REQUEST_METHOD'] === 'POST') {
    // 1) protection CSRF : comparaison du jeton recu avec celui de la session
    if (!hash_equals($csrf, (string) ($_POST['_csrf'] ?? ''))) {
        http_response_code(419);
        exit('Formulaire expire, recharge la page.');
    }
    // 2) recuperer et valider
    $valeurs['pseudo'] = champ($_POST, 'pseudo') ?? '';
    $valeurs['email'] = mb_strtolower(champ($_POST, 'email') ?? '');
    $valeurs['message'] = champ($_POST, 'message') ?? '';
    $valeurs['newsletter'] = isset($_POST['newsletter']) ? '1' : '';

    if (mb_strlen($valeurs['pseudo']) < 3) { $erreurs['pseudo'] = 'Au moins 3 caracteres.'; }
    if ($valeurs['email'] === '' || filter_var($valeurs['email'], FILTER_VALIDATE_EMAIL) === false) {
        $erreurs['email'] = 'Adresse email invalide.';
    }
    if (mb_strlen($valeurs['message']) > 1000) { $erreurs['message'] = '1000 caracteres maximum.'; }

    if ($erreurs === []) {
        // 3) ecrire (requete preparee : voir chapitre base de donnees)
        $stmt = $pdo->prepare('INSERT INTO inscriptions (pseudo, email, message, newsletter, cree_le) VALUES (?, ?, ?, ?, NOW())');
        $stmt->execute([$valeurs['pseudo'], $valeurs['email'], $valeurs['message'], $valeurs['newsletter'] === '1' ? 1 : 0]);
        $_SESSION['flash']['succes'] = 'Merci, c\'est enregistre.';
        rediriger('/inscription');           // PRG : un rafraichissement ne re-soumet pas
    }
    // sinon : on retombe dans le rendu avec $valeurs et $erreurs remplis
}

function e(?string $v): string { return htmlspecialchars((string) $v, ENT_QUOTES | ENT_SUBSTITUTE, 'UTF-8'); }
?>
<?php if (isset($_SESSION['flash']['succes'])): ?>
  <p class="succes"><?= e($_SESSION['flash']['succes']) ?></p>
  <?php unset($_SESSION['flash']['succes']); ?>
<?php endif; ?>

<form method="post" action="/inscription" novalidate enctype="application/x-www-form-urlencoded">
  <input type="hidden" name="_csrf" value="<?= e($csrf) ?>">

  <label for="pseudo">Pseudo</label>
  <input id="pseudo" name="pseudo" required minlength="3" maxlength="40"
         value="<?= e($valeurs['pseudo']) ?>"
         <?php if (isset($erreurs['pseudo'])): ?>aria-invalid="true" aria-describedby="err-pseudo"<?php endif; ?>>
  <?php if (isset($erreurs['pseudo'])): ?><p id="err-pseudo" class="erreur"><?= e($erreurs['pseudo']) ?></p><?php endif; ?>

  <label for="email">Email</label>
  <input id="email" name="email" type="email" required value="<?= e($valeurs['email']) ?>">
  <?php if (isset($erreurs['email'])): ?><p class="erreur"><?= e($erreurs['email']) ?></p><?php endif; ?>

  <label for="message">Message <small>(<?= mb_strlen($valeurs['message']) ?>/1000)</small></label>
  <textarea id="message" name="message" rows="6"><?= e($valeurs['message']) ?></textarea>

  <label><input type="checkbox" name="newsletter" value="1" <?= $valeurs['newsletter'] ? 'checked' : '' ?>> Newsletter</label>

  <button type="submit">Envoyer</button>
</form>
\end{codephp}

\begin{enumerate}
  \item \textbf{Toujours} renvoyer les valeurs saisies (sinon l'utilisateur retape tout) et
        ne jamais r\'eafficher un mot de passe.
  \item \textbf{Toujours} \'echapper \`a la sortie~: \verb|htmlspecialchars| avec
        \verb|ENT\_QUOTES| et UTF-8. Une vue qui n'\'echappe pas est une faille XSS.
  \item \textbf{Toujours} faire un POST/REDIRECT/GET apr\`es un envoi r\'eussi.
  \item \textbf{Toujours} valider c\^ot\'e serveur, m\^eme si le HTML a \verb|required|~: le HTML se
        contourne en dix secondes avec DevTools ou \verb|curl|.
  \item \textbf{Jamais} faire confiance \`a \verb|\$\_FILES['type']|, \`a la longueur annonc\'ee, ni \`a
        un \verb|name|~: tout cela vient du client.
\end{enumerate}

\section{Sessions~: le stockage c\^ot\'e serveur}

\begin{codephp}
<?php
// php.ini (ou ini_set avant session_start, ou dans un pool php-fpm) : les reglages qui comptent
ini_set('session.use_strict_mode', '1');              // rejette un ID de session fourni par l'attaquant
ini_set('session.cookie_httponly', '1');              // pas accessible en JS
ini_set('session.cookie_secure', '1');                 // HTTPS uniquement
ini_set('session.cookie_samesite', 'Lax');
ini_set('session.use_only_cookies', '1');              // jamais d'ID dans l'URL
ini_set('session.gc_maxlifetime', '1800');             // 30 min d'inactivite
ini_set('session.sid_length', '48'); ini_set('session.sid_bits_per_character', '6');

session_start();                                       // cree ou reprend la session
$_SESSION['user_id'] = 7;                              // uniquement des donnees petites et serialisables
$uid = $_SESSION['user_id'] ?? null;
session_regenerate_id(true);                            // APRES toute connexion/escalade : anti-fixation
session_destroy();                                      // suppression cote serveur
setcookie('panier', '', time() - 3600, '/');             // effacer un cookie = le reemettre expire
$nom = session_name();                                   // 'PHPSESSID' par defaut ; renomme-le (minorite d'infos)

// Session en base (multiprobe : plusieurs serveurs php-fpm) : SessionHandlerInterface
final class SessionBdd implements SessionHandlerInterface
{
    public function __construct(private readonly PDO $pdo, private readonly int $duree = 1800) {}
    public function open(string $path, string $name): bool { return true; }
    public function close(): bool { return true; }
    public function read(string $id): string|false
    {
        $s = $this->pdo->prepare('SELECT data FROM sessions WHERE id = ? AND expire > NOW()');
        $s->execute([$id]);
        return (string) ($s->fetchColumn() ?: '');
    }
    public function write(string $id, string $data): int|bool
    {
        $s = $this->pdo->prepare('REPLACE INTO sessions (id, data, expire) VALUES (?, ?, NOW() + INTERVAL ? SECOND)');
        return $s->execute([$id, $data, $this->duree]);
    }
    public function destroy(string $id): int|bool { return $this->pdo->prepare('DELETE FROM sessions WHERE id = ?')->execute([$id]); }
    public function gc(int $max_lifetime): int|false { return $this->pdo->exec('DELETE FROM sessions WHERE expire < NOW()'); }
}
(new SessionHandler())->... ; // en pratique : session_set_save_handler(new SessionBdd($pdo), true);
\end{codephp}

\begin{note}
\textbf{Ce qu'on ne met jamais en session~:} un gros tableau (la session est s\'erialis\'ee et
relue \`a chaque requ\^ete), un secret non chiffr\'e, un r\^ole recalculable~; et surtout pas un jeton
d'API dans une session stock\'ee en base. Ce qu'on y met~: l'identifiant utilisateur, un
\verb|csrf|, un \'etat d'\'etapes, des flash.
\end{note}

\section{Authentification compl\`ete~: le parcours}

\begin{codephp}
<?php
declare(strict_types=1);

final class Authentification
{
    public function __construct(
        private readonly PDO $pdo,
        private readonly int $dureeVerrou = 900,
        private readonly int $essaisMax = 5,
    ) {}

    /** Inscription : verifie la force, hache, ecrit. Reponse identique que l'email existe ou non. */
    public function inscrire(string $email, string $motDePasse, string $pseudo): void
    {
        $this->verifierForce($motDePasse);
        $hash = password_hash($motDePasse, PASSWORD_ARGON2ID, [
            'memory_cost' => 1 << 15, 'time_cost' => 3, 'threads' => 2,   // ~0,3 s : assez lent pour gener, assez rapide pour toi
        ]);
        $s = $this->pdo->prepare(
            'INSERT INTO utilisateurs (email, mot_de_passe, pseudo, cree_le) VALUES (?, ?, ?, NOW())
             ON DUPLICATE KEY UPDATE id = id'     ;
        );
        $s->execute([mb_strtolower($email), $hash, $pseudo]);
    }

    public function connecter(string $email, string $motDePasse, ?string $ip = null): ?Utilisateur
    {
        $email = mb_strtolower(trim($email));
        $s = $this->pdo->prepare('SELECT id, email, mot_de_passe, pseudo, verrou_jusqua FROM utilisateurs WHERE email = ? LIMIT 1');
        $s->execute([$email]);
        $u = $s->fetch(PDO::FETCH_ASSOC);

        // 1) reponse au meme rythme et au meme message que si l'email etait inconnu
        if ($u === false) {
            password_verify($motDePasse, '$argon2id$v=19$m=65536,t=3,p=2$AAAAAAAAAAAAAAAAAAAAAA$AAAAAAAAAAAAAAAAAAAAAAAAAAAAAAAAAAAAAAAAAAA');
            throw new ErreurAuthentification('Identifiants incorrects.');
        }
        if ($u['verrou_jusqua'] !== null && strtotime((string) $u['verrou_jusqua']) > time()) {
            throw new ErreurAuthentification('Trop d\'essais. Reessaie plus tard.');
        }
        if (!password_verify($motDePasse, (string) $u['mot_de_passe'])) {
            $this->marquerEchec((int) $u['id'], $ip);
            throw new ErreurAuthentification('Identifiants incorrects.');
        }
        // 2) re-hachage si l'algorithme ou les couts ont change : la maintenance de securite offerte gracieusement
        if (password_needs_rehash((string) $u['mot_de_passe'], PASSWORD_ARGON2ID,
            ['memory_cost' => 1 << 15, 'time_cost' => 3, 'threads' => 2])) {
            $n = password_hash($motDePasse, PASSWORD_ARGON2ID, ['memory_cost' => 1 << 15, 'time_cost' => 3, 'threads' => 2]);
            $this->pdo->prepare('UPDATE utilisateurs SET mot_de_passe = ? WHERE id = ?')->execute([$n, $u['id']]);
        }
        // 3) ouverture de session securisee
        session_regenerate_id(true);
        $_SESSION = ['user_id' => (int) $u['id'], 'ip' => $ip, 'agent' => hash('sha256', (string) ($_SERVER['HTTP_USER_AGENT'] ?? '')), 'depuis' => time()];
        $this->pdo->prepare('UPDATE utilisateurs SET verrou_jusqua = NULL, tentatives = 0, derniere_connexion = NOW() WHERE id = ?')
                 ->execute([$u['id']]);
        return Utilisateur::deLigne($u);
    }

    public function deconnecter(): void
    {
        if (ini_get('session.use_cookies')) {
            $p = session_get_cookie_params();
            setcookie(session_name(), '', time() - 42000, $p['path'], $p['domain'], $p['secure'], $p['httponly']);
        }
        $_SESSION = [];
        session_destroy();
    }

    /** Reinitialisation : jeton aleatoire, HACHE en base, 15 min, usage unique, reponse neutre. */
    public function demanderReinitialisation(string $email): void
    {
        $s = $this->pdo->prepare('SELECT id FROM utilisateurs WHERE email = ?');
        $s->execute([mb_strtolower($email)]);
        $id = $s->fetchColumn();
        if ($id === false) { return; }                      // pas de fuite d'information : meme message
        $brut = bin2hex(random_bytes(32));
        $this->pdo->prepare('INSERT INTO reinitialisations (utilisateur_id, jeton_hash, expire_le) VALUES (?, ?, NOW() + INTERVAL 15 MINUTE)')
                  ->execute([$id, hash('sha256', $brut)]);
        $this->mailer->envoyer($email, 'Reinitialisation', "Lien : " . url('/mdp-oublier') . '?j=' . urlencode($brut));
    }

    public function appliquerReinitialisation(string $jetonRecu, string $nouveauMdp): void
    {
        $this->verifierForce($nouveauMdp);
        $s = $this->pdo->prepare(
            'SELECT utilisateur_id FROM reinitialisations
             WHERE jeton_hash = ? AND expire_le > NOW() AND utilise_le IS NULL LIMIT 1'
        );
        $s->execute([hash('sha256', $jetonRecu)]);
        $id = $s->fetchColumn();
        if ($id === false) { throw new ErreurAuthentification('Lien invalide ou expire.'); }
        $this->pdo->prepare('UPDATE utilisateurs SET mot_de_passe = ? WHERE id = ?')
                  ->execute([password_hash($nouveauMdp, PASSWORD_ARGON2ID), $id]);
        $this->pdo->prepare('UPDATE reinitialisations SET utilise_le = NOW() WHERE utilisateur_id = ? AND utilise_le IS NULL')->execute([$id]);
        $this->pdo->prepare('DELETE FROM sessions WHERE utilisateur_id = ?')->execute([$id]);   // deconnecter partout
    }

    private function verifierForce(string $mdp): void
    {
        if (mb_strlen($mdp) < 12) { throw new ErreurValidation('Mot de passe trop court (12 caracteres minimum).'); }
        // (bonus) comparer aux fuites connues via l'API Have I Been Pwned range, sans envoyer le mot de passe
    }
    private function marquerEchec(int $id, ?string $ip): void { /* incrementer tentatives, poser verrou */ }
}
\end{codephp}

\subsection{Le middleware d'authentification, et les autorisations}

\begin{codephp}
<?php
// Un controleur ne verifie pas qui il est : un middleware le fait, une fois, pour tout.
final class ExigerConnexion
{
    public function __construct(private readonly PDO $pdo) {}
    public function __invoke(callable $suite): Response
    {
        $id = $_SESSION['user_id'] ?? null;
        if ($id === null) { return new Response(401, ['www-authenticate' => 'session'], 'Connexion requise'); }
        // Lier la session a l'empreinte du navigateur et a l'IP : casse la session volee depuis un autre poste
        if (!hash_equals($_SESSION['agent'] ?? '', hash('sha256', (string) ($_SERVER['HTTP_USER_AGENT'] ?? '')))) {
            $this->deconnecter();
            return new Response(401, [], 'Session invalide');
        }
        $u = Utilisateur::trouver($this->pdo, (int) $id) ?? (function (): never { $this->deconnecter(); throw new RuntimeException(); })();
        $GLOBALS['utilisateur'] = $u;             // dans un vrai framework : attribut de requete
        return $suite();
    }
    private function deconnecter(): void { (new Authentification($this->pdo))->deconnecter(); }
}

// Autorisation : centralisee, testable, et qui renvoie 404 quand l'existence meme est sensible
final class Politiques
{
    public static function peutVoir(Utilisateur $u, Article $a): bool
    {
        return $a->estPublie() || $u->id() === $a->auteur()->id() || $u->role()->peut(Role::Redacteur);
    }
    public static function peutEditer(Utilisateur $u, Article $a): bool
    {
        return $u->id() === $a->auteur()->id() || $u->role()->peut(Role::Admin);
    }
}
\end{codephp}

\section{Upload de fichiers~: la version non-na\"ive}

\begin{codephp}
<?php
/**
 * @param array{name:string,type:string,tmp_name:string,error:int,size:int} $f  $_FILES['avatar']
 */
function traiterUpload(array $f, string $dossier, int $tailleMax = 5_242_880): string
{
    if ($f['error'] !== UPLOAD_ERR_OK) {
        throw new ErreurUpload(match ($f['error']) {
            UPLOAD_ERR_INI_SIZE, UPLOAD_ERR_FORM_SIZE => 'Fichier trop volumineux.',
            UPLOAD_ERR_PARTIAL => 'Envoi interrompu.',
            UPLOAD_ERR_NO_FILE => 'Aucun fichier recu.',
            default => 'Erreur d\'envoi (' . $f['error'] . ').',
        });
    }
    if ($f['size'] > $tailleMax) { throw new ErreurUpload('Taille maximale depassee.'); }

    // 1) verifier que c'est REELLEMENT une image (le MIME annonce par le client ne vaut rien)
    $info = @getimagesize($f['tmp_name']);
    if ($info === false) { throw new ErreurUpload('Ce n\'est pas une image.'); }
    $types = [IMAGETYPE_PNG => 'png', IMAGETYPE_JPEG => 'jpg', IMAGETYPE_WEBP => 'webp', IMAGETYPE_GIF => 'gif'];
    $ext = $types[$info[2]] ?? throw new ErreurUpload('Format non accepte.');

    // 2) nom genere : jamais le nom du client dans un chemin (traversee de repertoire)
    $nom = sprintf('%s-%s.%s', date('Ymd'), bin2hex(random_bytes(8)), $ext);

    // 3) hors racine web,permissions 0640, proprietaire du service ; on reencode pour neutraliser
    //    toute charge utile cachee (EXIF, script) : passer par GD ou ImageMagick
    $destination = rtrim($dossier, '/') . '/' . $nom;
    if (!is_dir($dossier)) { mkdir($dossier, 0o750, true); }
    if (!move_uploaded_file($f['tmp_name'], $destination)) {      // LA seule fonction qui prouve
        throw new ErreurUpload('Enregistrement impossible.');                   // l'origine PHP de l'upload
    }
    reencoder($destination, $ext);                                  // redimensionne + retire EXIF
    chmod($destination, 0o640);
    return '/uploads/' . $nom;
}
// En nginx/apache : le dossier d'upload ne doit JAMAIS executer de PHP :
//   php_admin_flag engine Off / ou alias en `location ^~ /uploads/ { autoindex off; }`
\end{codephp}

\subsection{En-t\^etes, cookies, emails~: les trois finitions}

\begin{codephp}
<?php
// Cookies (le meme contrat que vu en JS)
setcookie('theme', 'sombre', [
    'expires' => time() + 31536000, 'path' => '/', 'domain' => '',
    'secure' => true, 'httponly' => true, 'samesite' => 'Lax', 'partitioned' => false,
]);
setrawcookie('brut', rawurlencode($valeur), ['expires' => 0, 'path' => '/']);

// En-tetes de securite, a poser globalement (php-fpm : dans le pool, ou dans le premier include)
header_remove('x-powered-by'); header_remove('server');
header('x-content-type-options: nosniff');
header('referrer-policy: strict-origin-when-cross-origin');
header('x-frame-options: DENY');
header('permissions-policy: camera=(), microphone=(), geolocation=()');
header('content-security-policy: default-src \'self\'; img-src \'self\' data:; style-src \'self\' \'unsafe-inline\'');
header('strict-transport-security: max-age=31536000; includeSubDomains');

// Email : la fonction mail() de PHP est une source de problems (pas d'auth, pas de file,
// spam). Envoie via SMTP reel ou une API. Deux choix sains :
//   - Symfony Mailer : composer require symfony/mailer
//   - PHPMailer     : composer require phpmailer/phpmailer
use Symfony\Component\Mailer\Mailer;
use Symfony\Component\Mailer\Transport\Smtp\EsmtpTransport;
use Symfony\Component\Mime\Email;

$transport = new EsmtpTransport((string) getenv('SMTP_HOST'), (int) getenv('SMTP_PORT'), true);
$transport->setUsername((string) getenv('SMTP_USER'))->setPassword((string) getenv('SMTP_PASS'));
(new Mailer($transport))->send(
    (new Email())->from('nour@example.com')->to($email)->subject('Bienvenue')
        ->text("Salut {$pseudo},\n\nTon compte est pret.")
        ->html('<p>Salut <b>' . htmlspecialchars($pseudo) . '</b></p>')
        ->attachFromPath($cheminPdf, 'guide.pdf')
);
\end{codephp}

\begin{exercice}
\begin{enumerate}
  \item \'Ecris un \og ~formulaire g\'en\'erique~ \fg{}~: \verb|champ(type, nom, options)| qui g\'en\`ere le HTML,
        lit \verb|\$\_POST|, valide selon un tableau de r\`egles, et renvoie
        \verb|[valeurs, erreurs]|. Trois formulaires deviennent alors dix lignes chacun.
  \item Ajoute la protection \og ~honeypot~ \fg{} (un champ cach\'e que seul un robot remplit) et un
        d\'elai minimal entre affichage et envoi~; bloque les robots sans CAPTCHA.
  \item \'Ecris le module d'inscription/connexion complet ci-dessus sur ton blog~: avec verrouillage
        apr\`es cinq \'echecs, journalisation des connexions, et \og ~se souvenir de moi~ \fg{} via un jeton
        en base (jamais un cookie \og ~mot de passe hach\'e~ \fg{}).
  \item Ajoute un 2FA TOTP~: g\'en\`ere le secret, le code de v\'erification (\link{otphp}{https://github.com/Spomky-Labs/otphp}),
        et exige le code pour les r\^oles sensibles.
  \item Fais un upload multiple (jusqu'\`a 5 images)~: validit\'e, redimensionnement en trois
        tailles, nommage, suppression \`a la suppression de l'article, et tests sur des fichiers
        pi\`eges (\verb|photo.php|.jpg, image SVG, image de 40~MP, GIF anim\'e).
  \item \'Ecris une page \og ~mes sessions actives~ \fg{}~: liste des sessions de l'utilisateur (agent,
        IP, date), et bouton pour les fermer une par une~; c'est l'exercice qui fait comprendre
        le stockage de session.
\end{enumerate}
\end{exercice}

\begin{ressources}
\begin{itemize}
  \item Manuel --- Gestion des sessions --- \url{https://www.php.net/manual/en/book.session.php}
  \item Manuel --- Fonction \verb|password\_*| --- \url{https://www.php.net/manual/en/ref.password.php}
  \item PHP-FIG PSR-7 (messages HTTP) --- \url{https://www.php-fig.org/psr/psr-7/} --- PSR-15 (middlewares)
  \item \link{Symfony Mailer}{https://symfony.com/doc/current/mailer.html} --- \link{PHPMailer}{https://github.com/PHPMailer/PHPMailer}
  \item \link{OWASP~Cheat~Sheet~Authentication}{https://cheatsheetseries.owasp.org/cheatsheets/Authentication\_Cheat\_Sheet.html}
\end{itemize}
\end{ressources}

\index{PHP!formulaire}\index{PHP!session}\index{PHP!authentification}\index{CSRF}\index{upload}
